%%
%% Director-CenterNet: Multi-Region Viewport Prediction for Automated Esports Broadcasting
%% Target Journal: Expert Systems with Applications (Elsevier)
%% Template: elsarticle-template-num.tex (numbered citation style)
%%

\documentclass[preprint,12pt]{elsarticle}

%% Use the option review to obtain double line spacing
%% \documentclass[authoryear,preprint,review,12pt]{elsarticle}

%% Use the options 1p,twocolumn; 3p; 3p,twocolumn; 5p; or 5p,twocolumn
%% for a journal layout:
%% \documentclass[final,1p,times]{elsarticle}
%% \documentclass[final,3p,times]{elsarticle}
%% \documentclass[final,5p,times]{elsarticle}

%% The amssymb package provides various useful mathematical symbols
\usepackage{amssymb}
%% The amsmath package provides various useful equation environments.
\usepackage{amsmath}
%% Extended theorem environments
%% \usepackage{amsthm}

%% Line numbers for review
%% \usepackage{lineno}

%% Additional packages
\usepackage{booktabs}
\usepackage{multirow}
\usepackage{algorithm}
\usepackage{algorithmic}
\usepackage{xcolor}
\usepackage{subcaption}
\usepackage{mathtools}

%% Convenience macros
\newcommand{\R}{\mathbb{R}}
\newcommand{\todo}[1]{\textcolor{red}{[\textbf{TODO}: #1]}}

\journal{Expert Systems with Applications}

\begin{document}

\begin{frontmatter}

%% Title
\title{Director-CenterNet: Multi-Region Viewport Prediction\\
       for Automated Esports Broadcasting}

%% Authors
%% use the corref command within \author for corresponding author footnote
\author[gist]{Cheong-mok Bae}
\ead{cmbae0307@gm.gist.ac.kr}

\author[gist]{Ho-Taek Joo}
\ead{hotaek87@gm.gist.ac.kr}

\author[aigist]{SungHa Lee}
\ead{shlee0414@gm.gist.ac.kr}

\author[gist,aigist]{Kyung-Joong Kim\corref{cor1}}
\ead{kjkim@gist.ac.kr}

\cortext[cor1]{Corresponding author}

\affiliation[gist]{organization={School of Integrated Technology},
             addressline={123 Cheomdan-gwagi-ro, Buk-gu},
             city={Gwangju},
             postcode={61005},
             country={South Korea}}

\affiliation[aigist]{organization={AI Graduate School},
             addressline={123 Cheomdan-gwagi-ro, Buk-gu},
             city={Gwangju},
             postcode={61005},
             country={South Korea}}

%% Abstract
\begin{abstract}
Automated esports broadcasting requires not only identifying the most 
engaging game region at each frame, but also simultaneously managing multiple 
camera feeds — a main screen and several picture-in-picture (PIP) channels — 
so that concurrent strategic events are not missed. Existing single-viewport 
observers based on behavior cloning from human observational data succeed at 
imitating individual camera placements but suffer from camera thrashing and 
cannot model the spatial resource allocation problem inherent in multi-screen 
broadcasting. In this paper, we formulate \emph{Multi-Region Viewport 
Prediction} (MRVP) as a structured prediction problem over a tile-aligned, 
multi-channel game-state tensor, and propose \emph{Director-CenterNet}, a 
heatmap-based multi-peak detection framework that jointly predicts a main 
viewport and up to $K$ auxiliary PIP viewports. Building on the Kernel-Based 
Region Scoring (KBRS) saliency prior — which encodes density, mixture, and 
centeredness cues into a fixed-kernel score field — we introduce four 
complementary training objectives: (i) a \emph{Human Consensus Match} loss 
anchoring the main query to the region of common interest among multiple human 
observers; (ii) a \emph{KBRS Coverage} loss directing sub-queries to salient 
regions not captured by the main viewport; (iii) a \emph{Spatial Repulsion} 
loss penalising redundant viewport overlap; and (iv) a \emph{Trajectory 
Smoothness} loss suppressing temporal camera thrashing. Experiments on 
StarCraft: Brood War replays processed through the Observer preprocessing 
pipeline demonstrate that Director-CenterNet outperforms all single-viewport 
baselines on Intersection Ratio and Intersection@\{0.3, 0.5\}, while 
producing temporally stable and spatially diverse viewport trajectories. 
Qualitative analysis confirms that the multi-region framework captures 
concurrent strategic events that single-viewport observers necessarily miss.
\end{abstract}

%% Graphical abstract
\begin{graphicalabstract}
%% \includegraphics{graphical_abstract}
\todo{Insert graphical abstract: pipeline diagram showing game-state tensor
$\rightarrow$ Director-CenterNet $\rightarrow$ Main + PIP viewports overlaid
on minimap.}
\end{graphicalabstract}

%% Research highlights
\begin{highlights}
\item We formulate Multi-Region Viewport Prediction (MRVP) as a structured
      prediction problem over multi-channel game-state tensors.
\item Director-CenterNet employs CenterNet heatmaps to jointly predict a main
      viewport and multiple PIP viewports in a single forward pass.
\item A KBRS Coverage loss allocates PIP viewports to salient uncovered regions
      without additional human annotations.
\item Spatial repulsion and trajectory smoothness losses eliminate redundant
      overlap and camera thrashing, respectively.
\item Director-CenterNet outperforms all single-viewport baselines in viewport
      overlap and temporal stability on StarCraft benchmarks.
\end{highlights}

%% Keywords
\begin{keyword}
Esports \sep Automatic observer \sep Multi-viewport prediction \sep
CenterNet \sep Saliency-guided learning \sep Spatial resource allocation \sep
StarCraft \sep Behavior cloning
\end{keyword}

\end{frontmatter}

%% \linenumbers   %% uncomment for review line numbers

%%====================================================================
\section{Introduction}
\label{sec:intro}
%%====================================================================

Human game observers in electronic sports (Esports) fulfill a dual role: 
they must simultaneously identify the most strategically significant region 
of the map \emph{and} coordinate multiple camera feeds to ensure that 
concurrent events are not missed \cite{joo2023learning}. In professional 
Esports productions, a team of observers typically manages one main broadcast 
camera together with several auxiliary picture-in-picture (PIP) channels, each 
covering a distinct map region. Automating this workflow is challenging for 
three reasons: (1) the number of concurrent high-priority events is dynamic; 
(2) the spatial relationships among events evolve rapidly; and (3) the 
subjective notion of ``broadcast saliency'' cannot be reduced to a single, 
well-defined objective \cite{bae2025observer}.

Prior work on automatic Esports observers has focused on the 
\emph{single-viewport} setting, where the goal is to select one rectangular 
window of fixed size from the game map at each frame. Rule-based approaches 
\cite{mattsson2015automatic} define priority scores over hand-crafted events 
and switch the camera accordingly. Learning-based methods 
\cite{joo2023learning} instead collect observational data from human 
spectators and train a detector-style model with Region of Common Interest 
(ROCI) loss to imitate consensus human camera behaviour. The Kernel-Based 
Region Scoring (KBRS) auxiliary loss \cite{anonymous2025kbrs} further improved 
alignment with aggregated human attention by injecting fixed-kernel saliency 
priors (density, mixture, centeredness) into the detection backbone at 
training time.

Despite these advances, two fundamental limitations remain open. First, 
\textbf{single-viewport observers cannot cover simultaneous events}: when a 
battle occurs at one map corner while a covert resource harassment unfolds 
elsewhere, no single viewport can serve both. Second, 
\textbf{behaviour-cloned single-viewport models exhibit camera thrashing}: 
because the model is trained to match whichever event the majority of human 
observers happened to watch, it oscillates between competing salient regions 
across frames, producing broadcast footage that is visually unstable.

We address both limitations by reformulating automatic observing as 
\textbf{Multi-Region Viewport Prediction (MRVP)}: at each frame, the system 
must produce a \emph{set} of $K$ non-redundant viewports — one main viewport 
aligned with human consensus attention and $K-1$ auxiliary PIP viewports 
covering other salient regions not captured by the main screen. This 
formulation naturally models the multi-screen broadcasting workflow while 
suppressing thrashing through spatial repulsion constraints.

To solve MRVP we propose \textbf{Director-CenterNet}, built on CenterNet 
\cite{zhou2019objects}. CenterNet's dense heatmap representation is 
well-suited to MRVP because (i) it simultaneously scores all candidate 
viewport centres, enabling multi-peak extraction without iterative proposal 
generation; (ii) the heatmap can be directly regularised by the KBRS score 
field during training; and (iii) local-maximum extraction provides a 
principled mechanism for dynamic PIP management — suppressing a PIP channel 
when no secondary salient region is detected.

Our contributions are as follows.

\begin{enumerate}
  \item We formalise Multi-Region Viewport Prediction (MRVP) and introduce 
        new evaluation metrics for multi-viewport coverage and temporal 
        stability (Section~\ref{sec:problem}).
  \item We propose Director-CenterNet, a CenterNet-based multi-peak heatmap 
        observer that jointly predicts main and PIP viewports 
        (Section~\ref{sec:method}).
  \item We introduce a KBRS Coverage loss that explicitly directs PIP queries 
        to high-KBRS-score regions not covered by the main viewport, enabling 
        principled spatial resource allocation without additional human 
        annotations (Section~\ref{subsec:cov_loss}).
  \item We demonstrate superior viewport overlap and temporal stability 
        compared to all single-viewport baselines on StarCraft: Brood War 
        benchmarks (Section~\ref{sec:experiments}).
\end{enumerate}

%%====================================================================
\section{Background and Related Work}
\label{sec:related}
%%====================================================================

\subsection{Single-Viewport Automatic Observers}

Automatic game observation has been studied primarily in real-time strategy 
(RTS) games, where the viewport selection problem is most challenging due to 
large map size and high event concurrency \cite{ontanon2013survey}. 
Rule-based observers such as AIIDE and SSCAIT \cite{mattsson2015automatic} 
define priority scores over game events and teleport the camera to the 
highest-priority event at each decision step. While interpretable, these 
systems require manual event taxonomies and are brittle under novel strategies.

Joo et al. \cite{joo2023learning} proposed the first event-free learning-based 
observer for StarCraft. They framed viewport selection as an instance 
segmentation problem, treating the human observer's viewport as the object to 
detect. Using Mask R-CNN \cite{he2020mask} and a Region of Common Interest 
(ROCI) loss that amplifies the overlapping area among multiple human observers, 
their method outperforms both rule-based baselines and behaviour cloning. 
Anonymous et al. \cite{anonymous2025kbrs} further improved upon 
\cite{joo2023learning} by introducing Kernel-Based Region Scoring (KBRS) as 
a training-time auxiliary loss. KBRS computes three fixed-kernel score maps — 
density ($S^{\text{den}}$), mixture ($S^{\text{mix}}$), and centeredness 
($S^{\text{ctr}}$) — on backbone FPN features, improving intersection ratio 
by 4.31\% relative without altering inference-time outputs.

The Observer framework \cite{bae2025observer} provides an open-source, 
reproducible pipeline for data collection, preprocessing, and model training 
on StarCraft: Brood War replays, which we use in this work.

\subsection{Object Detection Architectures}

Our work builds on two complementary detection paradigms. 
\textbf{CenterNet} \cite{zhou2019objects} frames detection as keypoint 
heatmap regression: the network predicts a Gaussian heatmap peaked at each 
object's centre, with offset and size regression heads for bounding box 
recovery. Multi-peak extraction by local-maximum search provides a natural 
mechanism for detecting a variable number of objects per image. 
\textbf{DETR} \cite{carion2020end} and RT-DETR \cite{lv2023detrs} reframe 
detection as set prediction via Hungarian matching, eliminating 
non-maximum suppression. We choose CenterNet over transformer-based detectors 
for MRVP because its dense heatmap output can be directly compared to the KBRS 
score field and supports dynamic PIP count without modifying the query set.

\subsection{Saliency and Multi-Region Detection}

The problem of distributing a limited number of viewports over a dynamic scene 
relates to multi-salient region detection \cite{shi2016hierarchical}, active 
perception \cite{chen2011active}, and sensor placement 
\cite{krause2008near}. Classical saliency models incorporate hand-designed 
priors — including centre bias and background priors — to guide attentional 
selection \cite{itti1998model, wei2012geodesic, tatler2007central}. Our work 
differs in that salient regions are defined by aggregated human attention over 
a structured multi-channel game-state tensor rather than by low-level visual 
features in natural images.

\subsection{Temporal Stability in Video-Based Observers}

Camera thrashing — rapid oscillation between competing salient regions — is a 
known failure mode in learning-based observers. Trajectory smoothness losses 
penalising inter-frame displacement have been applied in video object detection 
\cite{feichtenhofer2017detect} and automated cinematography 
\cite{gandhi2014video} to reduce temporal inconsistency. We adapt this 
approach to the MRVP setting by penalising the L2 displacement of predicted 
viewport centres across consecutive frames.

%%====================================================================
\section{Problem Formulation}
\label{sec:problem}
%%====================================================================

\subsection{Input Representation}

Following \cite{joo2023learning, bae2025observer, anonymous2025kbrs}, at 
each game frame $t$ we construct a multi-channel state tensor 
$\mathbf{X}_t \in \mathbb{R}^{C \times H \times W}$ on a tile grid of size 
$H \times W$. Channels encode semantically distinct entity categories for 
each player (workers, ground units, air units, buildings), terrain elevation, 
combined vision, and resource locations. Each tile corresponds to 
$32 \times 32$ pixels in the game world.

To incorporate short-term temporal context, we stack frames at a fixed 
stride $\Delta$:
\begin{equation}
  \mathbf{X}^{\text{stk}}_t 
    = \text{concat}\bigl[\mathbf{X}_t,\,\mathbf{X}_{t-\Delta},\,
      \ldots,\,\mathbf{X}_{t-W\Delta}\bigr]
    \in \mathbb{R}^{(W+1)C \times H \times W}.
  \label{eq:stacking}
\end{equation}
Unless otherwise stated we use $\Delta = 8$ and $W = 3$, yielding a 
four-frame stack.

\subsection{Human Supervision and Aggregation}

For each frame $t$ and human observer $u \in \{1,\ldots,U\}$, we record a 
viewport $V^{(u)}_t$ specified by its top-left tile coordinate 
$\mathbf{p}^{(u)}_t \in \mathcal{C}$, where
\begin{equation}
  \mathcal{C} = \{(i,j) \mid 0 \le i \le H - h,\; 0 \le j \le W - w\}
  \label{eq:feasible}
\end{equation}
is the set of feasible top-left positions for a viewport of fixed tile size 
$(h, w)$. We aggregate human viewports into a hard union set
\begin{equation}
  \mathcal{H}_t = \bigcup_{u=1}^{U} V^{(u)}_t
  \label{eq:union}
\end{equation}
and a soft crowd coverage map
\begin{equation}
  a_t(\mathbf{q}) = \frac{1}{U} \sum_{u=1}^{U} 
    \mathbf{1}\!\bigl[\mathbf{q} \in V^{(u)}_t\bigr],
  \label{eq:coverage}
\end{equation}
which measures the fraction of observers whose viewport covers tile 
$\mathbf{q}$.

\subsection{Multi-Region Viewport Prediction (MRVP)}

\textbf{Definition.} Given a game-state tensor $\mathbf{X}^{\text{stk}}_t$, 
predict a set of $K$ viewport positions 
$\{\hat{\mathbf{p}}^{(k)}_t\}_{k=1}^{K} \subset \mathcal{C}$ such that:
\begin{enumerate}
  \item $\hat{\mathbf{p}}^{(1)}_t$ (main viewport) maximally overlaps the 
        human consensus region $\mathcal{H}_t$;
  \item $\hat{\mathbf{p}}^{(k)}_t$ for $k \ge 2$ (PIP viewports) cover 
        salient regions not covered by the main viewport;
  \item Predicted viewports are spatially diverse (non-redundant);
  \item Predicted viewport positions vary smoothly across consecutive frames.
\end{enumerate}

This formulation generalises single-viewport prediction (the special case 
$K=1$) while connecting MRVP to the multi-salient region detection literature.

\subsection{Evaluation Metrics}

\paragraph{Intersection Ratio (IR).}
\begin{equation}
  \text{IR}_t = \frac{\bigl|V(\hat{\mathbf{p}}^{(1)}_t) \cap \mathcal{H}_t\bigr|}
                     {\bigl|V(\hat{\mathbf{p}}^{(1)}_t)\bigr|}
  \label{eq:ir}
\end{equation}

\paragraph{Intersection@$\delta$.}
\begin{equation}
  \text{I@}\delta_t = \mathbf{1}[\text{IR}_t \ge \delta],\quad
  \delta \in \{0,\, 0.3,\, 0.5\}.
  \label{eq:iat}
\end{equation}

\paragraph{Multi-Viewport Coverage (MVC).}
The fraction of $\mathcal{H}_t$ covered by the union of all $K$ predicted 
viewports:
\begin{equation}
  \text{MVC}_t = \frac{\displaystyle\left|\Bigl(\bigcup_{k=1}^{K} 
    V(\hat{\mathbf{p}}^{(k)}_t)\Bigr) \cap \mathcal{H}_t\right|}
    {|\mathcal{H}_t|}.
  \label{eq:mvc}
\end{equation}

\paragraph{Viewport Displacement (VD).}
Average L2 displacement of the main viewport centre across consecutive 
frames, measuring temporal stability:
\begin{equation}
  \text{VD} = \frac{1}{T-1}\sum_{t=1}^{T-1}
    \bigl\|\bar{\mathbf{p}}^{(1)}_t - \bar{\mathbf{p}}^{(1)}_{t+1}\bigr\|_2,
  \label{eq:vd}
\end{equation}
where $\bar{\mathbf{p}}^{(k)}_t$ denotes the centre of the $k$-th predicted 
viewport at frame $t$.

%%====================================================================
\section{Director-CenterNet}
\label{sec:method}
%%====================================================================

\subsection{Architecture Overview}

Director-CenterNet is a heatmap-based multi-peak detector that extends 
CenterNet \cite{zhou2019objects} to the MRVP problem. The architecture 
comprises three components: (1) a ResNet-FPN backbone encoder that transforms 
the stacked game-state tensor into a spatial feature map; (2) CenterNet heads 
that regress a dense Gaussian heatmap of viewport centres together with offset 
and size maps; and (3) the KBRS auxiliary module that computes three 
fixed-kernel score maps used exclusively during training. 
Fig.~\ref{fig:architecture} provides an overview.

\begin{figure}[t]
  \centering
  \todo{Insert architecture figure: Input tensor $\rightarrow$ ResNet-FPN 
  Backbone $\rightarrow$ CenterNet Heads (Heatmap / Offset / Size) + KBRS 
  Module (Density / Mixture / Centeredness) $\rightarrow$ Multi-peak 
  extraction $\rightarrow$ Main + PIP viewports.}
  \caption{Overview of Director-CenterNet. The stacked game-state tensor is 
  processed by a ResNet-FPN backbone. CenterNet heads output a dense heatmap 
  of viewport centres, along with offset and size regression maps. The KBRS 
  module computes three fixed-kernel score maps on FPN features and 
  contributes four auxiliary losses during training only; inference-time 
  predictions are produced solely by the vanilla CenterNet heads.}
  \label{fig:architecture}
\end{figure}

\subsection{CenterNet Backbone and Heads}

\paragraph{Input.} The stacked tensor 
$\mathbf{X}^{\text{stk}}_t \in \mathbb{R}^{(W+1)C \times H \times W}$ is 
fed to a ResNet-50 \cite{he2016deep} backbone with a Feature Pyramid Network 
(FPN) \cite{lin2017feature} to obtain a feature map 
$\mathbf{F}_t \in \mathbb{R}^{C_f \times H_s \times W_s}$ at stride $s$.

\paragraph{Heatmap head.} Two $3{\times}3$ convolutions followed by a 
$1{\times}1$ sigmoid layer produce the viewport-centre heatmap 
$\hat{\mathcal{Y}}_t \in [0,1]^{H_s \times W_s}$. The ground-truth heatmap 
is constructed by rendering a 2D Gaussian at each human-observed viewport 
centre $\bar{\mathbf{p}}^{(u)}_t$:
\begin{equation}
  \mathcal{Y}_t(\mathbf{q}) = 
    \max_u \exp\!\left(-\frac{\|\mathbf{q} - \bar{\mathbf{p}}^{(u)}_t\|^2}
                              {2\sigma^2}\right).
  \label{eq:gt_heatmap}
\end{equation}

\paragraph{Offset and size heads.} Following \cite{zhou2019objects}, 
a sub-pixel offset map $\hat{\mathbf{O}}_t \in \mathbb{R}^{2 \times H_s 
\times W_s}$ and a size map $\hat{\mathbf{Z}}_t \in \mathbb{R}^{2 \times 
H_s \times W_s}$ recover accurate viewport bounding boxes from detected 
centres.

\paragraph{Multi-peak extraction at inference.} A $3{\times}3$ max-pooling 
is applied to $\hat{\mathcal{Y}}_t$ to suppress non-maxima. The top-$K$ 
peaks with score above a confidence threshold $\tau$ are extracted as 
viewport centres. The main viewport corresponds to the highest-scoring peak; 
PIP viewports correspond to subsequent peaks. If fewer than $K$ peaks exceed 
$\tau$, the corresponding PIP channels are deactivated (dynamic PIP 
management).

\subsection{KBRS Saliency-Prior Score Field}

We retain the KBRS module from \cite{anonymous2025kbrs}. Given the FPN 
feature map $\mathbf{F}_t$, three fixed-kernel score maps are computed on the 
KBRS spatial grid of size $(H_s, W_s)$.

\paragraph{Density.} Let $\mathbf{K}_1 \in \mathbb{R}^{k_h \times k_w}$ be 
an all-ones window kernel. The density map accumulates local unit presence:
\begin{equation}
  S^{\text{den}}_t = \mathbf{K}_1 * 
    \!\left(\sum_{c \in \mathcal{G}} \mathbf{F}_t[c, \cdot, \cdot]\right),
  \label{eq:density}
\end{equation}
where $\mathcal{G}$ is the union of all entity channel groups.

\paragraph{Mixture.} Let $\mathcal{G}_1$ and $\mathcal{G}_2$ denote channel 
groups for player~1 and player~2 entities. Windowed group activations 
$u_t = \mathbf{K}_1 * \sum_{c \in \mathcal{G}_1}\mathbf{F}_t[c,\cdot,\cdot]$ 
and $v_t = \mathbf{K}_1 * \sum_{c \in \mathcal{G}_2}\mathbf{F}_t[c,\cdot,\cdot]$ 
form the local mixture proportion $\rho_t = u_t / (u_t + v_t + \varepsilon)$, 
and the mixture score is:
\begin{equation}
  S^{\text{mix}}_t = \bigl(4\rho_t(1-\rho_t)\bigr)^{\eta},
  \label{eq:mixture}
\end{equation}
with sharpness exponent $\eta \ge 1$.

\paragraph{Centeredness.} Let $\mathbf{K}_{\text{ctr}}$ be a fixed 
centre-biased Gaussian kernel:
\begin{equation}
  S^{\text{ctr}}_t = \mathbf{K}_{\text{ctr}} *
    \!\left(\sum_{c \in \mathcal{G}} \mathbf{F}_t[c, \cdot, \cdot]\right).
  \label{eq:centeredness}
\end{equation}

\paragraph{Combined score field.}
\begin{equation}
  S_t = \alpha S^{\text{den}}_t + \beta S^{\text{mix}}_t 
             + \gamma S^{\text{ctr}}_t,
  \label{eq:kbrs_combined}
\end{equation}
where $\alpha, \beta, \gamma \ge 0$ are hyperparameters.

\subsection{Multi-Objective Training Losses}
\label{subsec:losses}

Director-CenterNet is trained with four complementary objectives beyond 
the standard CenterNet detection losses.

\subsubsection{Human Consensus Match Loss $\mathcal{L}_{\text{hcm}}$}

The main viewport should be anchored to the region where the majority of 
human observers agreed to watch. The consensus centre is the centre of mass 
of the soft coverage map $a_t$ (Eq.~\ref{eq:coverage}):
\begin{equation}
  \mathbf{p}^*_t = \frac{\sum_{\mathbf{q}} a_t(\mathbf{q})\,\mathbf{q}}
                         {\sum_{\mathbf{q}} a_t(\mathbf{q})}.
  \label{eq:consensus_center}
\end{equation}
The ground-truth heatmap $\mathcal{Y}_t$ (Eq.~\ref{eq:gt_heatmap}) is 
rendered at all human viewport centres, and the standard modified focal loss 
\cite{law2018cornernet} is applied:
\begin{equation}
  \mathcal{L}_{\text{hcm}} = -\frac{1}{N}\sum_{\mathbf{q}}
    \begin{cases}
      (1-\hat{\mathcal{Y}}_t(\mathbf{q}))^{\alpha_f} 
        \log \hat{\mathcal{Y}}_t(\mathbf{q}) 
        & \text{if } \mathcal{Y}_t(\mathbf{q})=1, \\[4pt]
      (1-\mathcal{Y}_t(\mathbf{q}))^{\beta_f} 
        \hat{\mathcal{Y}}_t(\mathbf{q})^{\alpha_f} 
        \log(1-\hat{\mathcal{Y}}_t(\mathbf{q})) 
        & \text{otherwise,}
    \end{cases}
  \label{eq:hcm_loss}
\end{equation}
where $N$ is the number of ground-truth peaks and $\alpha_f, \beta_f$ are 
focal loss hyperparameters.

\subsubsection{KBRS Coverage Loss $\mathcal{L}_{\text{cov}}$}
\label{subsec:cov_loss}

PIP viewports should cover high-KBRS-score regions not already covered by 
the main viewport. Let $\mathcal{M}^{(1)}_t$ be the binary mask of tiles 
covered by the main viewport. The residual KBRS score field is:
\begin{equation}
  \tilde{S}_t(\mathbf{q}) = S_t(\mathbf{q}) \cdot 
    \bigl(1 - \mathcal{M}^{(1)}_t(\mathbf{q})\bigr).
  \label{eq:residual_kbrs}
\end{equation}
A target heatmap $\mathcal{Y}^{\text{cov}}_t$ is rendered via 2D Gaussian 
rendering at the top-$(K-1)$ peaks of $\tilde{S}_t$. The coverage loss is 
the mean-squared error between the PIP portion of the predicted heatmap and 
this target:
\begin{equation}
  \mathcal{L}_{\text{cov}} = \frac{1}{H_s W_s}\sum_{\mathbf{q}}
    \Bigl(\hat{\mathcal{Y}}_t(\mathbf{q})\cdot
      (1-\mathcal{M}^{(1)}_t(\mathbf{q})) 
      - \mathcal{Y}^{\text{cov}}_t(\mathbf{q})\Bigr)^2.
  \label{eq:cov_loss}
\end{equation}

\subsubsection{Spatial Repulsion Loss $\mathcal{L}_{\text{rep}}$}

To prevent multiple predicted viewports from covering the same region, we 
penalise pairwise Intersection over Union (IoU) between predicted viewports:
\begin{equation}
  \mathcal{L}_{\text{rep}} = \sum_{k < l} 
    \text{IoU}\!\bigl(V(\hat{\mathbf{p}}^{(k)}_t),\,
                      V(\hat{\mathbf{p}}^{(l)}_t)\bigr).
  \label{eq:rep_loss}
\end{equation}
When only one salient peak is detected, $\mathcal{L}_{\text{rep}} = 0$ and 
PIP channels are effectively suppressed.

\subsubsection{Trajectory Smoothness Loss $\mathcal{L}_{\text{smooth}}$}

To suppress camera thrashing, we penalise large inter-frame displacements 
of the main viewport centre:
\begin{equation}
  \mathcal{L}_{\text{smooth}} = \frac{1}{T-1}\sum_{t=1}^{T-1}
    \bigl\|\bar{\mathbf{p}}^{(1)}_t - \bar{\mathbf{p}}^{(1)}_{t+1}\bigr\|_2^2.
  \label{eq:smooth_loss}
\end{equation}

\subsubsection{Joint Training Objective}

The total loss combines the standard CenterNet detection objective with the 
four multi-viewport losses:
\begin{equation}
  \mathcal{L} = \mathcal{L}_{\text{hcm}} 
              + \lambda_{\text{off}}\mathcal{L}_{\text{off}} 
              + \lambda_{\text{sz}}\mathcal{L}_{\text{size}} 
              + \lambda_{\text{cov}}\mathcal{L}_{\text{cov}} 
              + \lambda_{\text{rep}}\mathcal{L}_{\text{rep}} 
              + \lambda_{\text{sm}}\mathcal{L}_{\text{smooth}},
  \label{eq:total_loss}
\end{equation}
where $\mathcal{L}_{\text{off}}$ and $\mathcal{L}_{\text{size}}$ are the 
standard CenterNet offset and size regression losses \cite{zhou2019objects}.

Table~\ref{tab:loss_summary} summarises the role of each loss component.

\begin{table}[t]
  \centering
  \caption{Summary of Director-CenterNet training losses.}
  \label{tab:loss_summary}
  \begin{tabular}{llll}
    \toprule
    Loss & Symbol & Type & Purpose \\
    \midrule
    Human Consensus Match & $\mathcal{L}_{\text{hcm}}$ & Focal & 
      Anchor main viewport to human consensus \\
    Offset regression & $\mathcal{L}_{\text{off}}$ & L1 & 
      Sub-pixel centre accuracy \\
    Size regression & $\mathcal{L}_{\text{size}}$ & L1 & 
      Viewport size accuracy \\
    KBRS Coverage & $\mathcal{L}_{\text{cov}}$ & MSE & 
      Direct PIP to uncovered KBRS peaks \\
    Spatial Repulsion & $\mathcal{L}_{\text{rep}}$ & IoU penalty & 
      Prevent redundant viewport overlap \\
    Trajectory Smoothness & $\mathcal{L}_{\text{smooth}}$ & L2 & 
      Suppress camera thrashing \\
    \bottomrule
  \end{tabular}
\end{table}

%%====================================================================
\section{Experiments}
\label{sec:experiments}
%%====================================================================

\subsection{Dataset and Preprocessing}

We use the same StarCraft: Brood War dataset and Observer preprocessing 
pipeline as \cite{joo2023learning, anonymous2025kbrs}. Replay files were 
collected from online communities (Ygosu, BWReplays) with constraints on 
player APM ($>250$), race matchups (Z vs.~P, T vs.~P, Z vs.~T), competitive 
maps (Fighting Spirit, Jade, Othello), and starting locations. Human 
observational data were collected from $U=25$ participants (aged 20--35) who 
freely spectated replay files while viewport coordinates were recorded at each 
frame. A total of approximately 420{,}000 frames from 21 games were collected.

Game-state tensors follow the channelisation protocol of \cite{bae2025observer}: 
four entity channels per player (workers, ground units, air units, buildings), 
terrain elevation, combined vision, and resource channels. Temporal stacking 
uses $\Delta=8$, $W=3$ (four-frame stack).

\subsection{Experimental Setup}

\paragraph{Splits.} Three replay-level folds are used for cross-validation 
following \cite{joo2023learning, anonymous2025kbrs}. All main comparison 
results are averaged over three folds and three random seeds 
$\{123, 456, 789\}$.

\paragraph{Implementation.} Backbone: ResNet-50 pretrained on ImageNet with 
FPN. Gaussian rendering: $\sigma = 2$~tiles. KBRS weights: 
$(\alpha,\beta,\gamma) = (0.3,\, 3.0,\, 0.3)$, window size $12{\times}20$ 
tiles (matching viewport size). Loss weights: $\lambda_{\text{off}}=1$, 
$\lambda_{\text{sz}}=0.1$, $\lambda_{\text{cov}}=0.5$, 
$\lambda_{\text{rep}}=0.3$, $\lambda_{\text{sm}}=0.2$. Multi-peak: $K=3$ 
(one main + two PIP), confidence threshold $\tau=0.2$.

\paragraph{Optimiser.} SGD with momentum~0.9, weight decay $5{\times}10^{-4}$, 
learning rate $5{\times}10^{-3}$ with StepLR decay (step~3, factor~0.1), 
batch size~16, 30~epochs.

\subsection{Baselines}

We compare Director-CenterNet against the following methods.

\begin{itemize}
  \item \textbf{BC}: Behaviour cloning via MSE regression to human viewport 
        coordinates.
  \item \textbf{AIIDE}: Rule-based observer \cite{mattsson2015automatic}.
  \item \textbf{SSCAIT}: Rule-based observer with event priorities and timers 
        \cite{mattsson2015automatic}.
  \item \textbf{Mask R-CNN (1H)}: Single-viewport detection, one human label 
        \cite{joo2023learning}.
  \item \textbf{Mask R-CNN (5H)}: Single-viewport detection, five human labels 
        \cite{joo2023learning}.
  \item \textbf{Mask R-CNN + ROCI}: With Region of Common Interest loss 
        \cite{joo2023learning}.
  \item \textbf{Mask R-CNN + KBRS}: With Kernel-Based Region Scoring 
        \cite{anonymous2025kbrs}.
  \item \textbf{CenterNet (single)}: CenterNet with single-viewport training 
        only (no multi-viewport losses), serving as a backbone ablation.
\end{itemize}

\subsection{Main Results}

Table~\ref{tab:main_results} reports main-viewport performance (IR and 
Intersection@$\delta$) on the standard test set.

\begin{table}[t]
  \centering
  \caption{Main viewport performance on the standard test set (fixed race, 
  map, starting location). Results averaged over three folds $\times$ three 
  seeds. Best result in \textbf{bold}, second best \underline{underlined}. 
  \todo{Fill in experimental values.}}
  \label{tab:main_results}
  \begin{tabular}{lcccc}
    \toprule
    Method & @any $\uparrow$ & @0.3 $\uparrow$ & @0.5 $\uparrow$ 
           & IR $\uparrow$ \\
    \midrule
    BC                         & 0.061 & 0.059 & 0.046 & -- \\
    AIIDE                      & 0.524 & --    & --    & -- \\
    SSCAIT                     & 0.491 & --    & --    & -- \\
    Mask R-CNN (1H)            & 0.593 & --    & --    & -- \\
    Mask R-CNN (5H)            & 0.569 & --    & --    & -- \\
    Mask R-CNN + ROCI          & 0.588 & --    & --    & -- \\
    Mask R-CNN + KBRS          & \underline{0.824} & \underline{0.798} 
                               & \underline{0.770} & \underline{0.709} \\
    CenterNet (single)         & \todo{X} & \todo{X} & \todo{X} & \todo{X} \\
    \midrule
    \textbf{Director-CenterNet (Ours)} & \todo{X} & \todo{X} 
                               & \todo{X} & \todo{X} \\
    \bottomrule
  \end{tabular}
\end{table}

Table~\ref{tab:mvc_results} reports Multi-Viewport Coverage (MVC) and 
Viewport Displacement (VD).

\begin{table}[t]
  \centering
  \caption{Multi-viewport coverage and temporal stability. MVC is computed 
  against the union of all five human viewports $\mathcal{H}_t$. Lower VD 
  is better. \todo{Fill in experimental values.}}
  \label{tab:mvc_results}
  \begin{tabular}{lccc}
    \toprule
    Method & MVC $\uparrow$ & VD (tiles) $\downarrow$ & Avg.\ active PIP \\
    \midrule
    Mask R-CNN + KBRS (single)       & -- & -- & 0 \\
    CenterNet (single)               & -- & \todo{X} & 0 \\
    \textbf{Director-CenterNet}      & \todo{X} & \todo{X} & \todo{X} \\
    \bottomrule
  \end{tabular}
\end{table}

\subsection{Generalisation Tests}

Following \cite{joo2023learning, anonymous2025kbrs}, we evaluate 
generalisation to unseen races, starting locations, and maps.

\begin{table}[t]
  \centering
  \caption{Generalisation test: IR (4/4 game progression) for unseen 
  conditions. \todo{Fill in experimental values.}}
  \label{tab:generalisation}
  \begin{tabular}{lccc}
    \toprule
    Method & Unseen Race & Unseen Location & Unseen Map \\
    \midrule
    Mask R-CNN + ROCI   & 0.615 & 0.559 & 0.581 \\
    Mask R-CNN + KBRS   & \todo{X} & \todo{X} & \todo{X} \\
    \textbf{Director-CenterNet} & \todo{X} & \todo{X} & \todo{X} \\
    \bottomrule
  \end{tabular}
\end{table}

\subsection{Ablation Study}

Table~\ref{tab:ablation} ablates each loss component to quantify individual 
contributions on fold~1, seed~123.

\begin{table}[t]
  \centering
  \caption{Ablation study: impact of each loss component 
  (fold 1, seed 123, 4/4 game progression). 
  \checkmark~= included. \todo{Fill in experimental values.}}
  \label{tab:ablation}
  \begin{tabular}{cccccccc}
    \toprule
    $\mathcal{L}_{\text{hcm}}$ & 
    $\mathcal{L}_{\text{cov}}$ & 
    $\mathcal{L}_{\text{rep}}$ & 
    $\mathcal{L}_{\text{smooth}}$ & 
    IR $\uparrow$ & MVC $\uparrow$ & VD $\downarrow$ \\
    \midrule
    \checkmark &            &            &            & \todo{X} & --       & --       \\
    \checkmark & \checkmark &            &            & \todo{X} & \todo{X} & --       \\
    \checkmark & \checkmark & \checkmark &            & \todo{X} & \todo{X} & --       \\
    \checkmark & \checkmark & \checkmark & \checkmark & \todo{X} & \todo{X} & \todo{X} \\
    \bottomrule
  \end{tabular}
\end{table}

\subsection{Qualitative Analysis}

\todo{Insert qualitative figures:
(1) Side-by-side: single-viewport vs.\ Director-CenterNet main + PIP during 
    concurrent game events (battle + resource harassment). 
(2) Heatmap visualisation: KBRS score field $S_t$, predicted heatmap 
    $\hat{\mathcal{Y}}_t$, residual field $\tilde{S}_t$ after main viewport 
    masking, and extracted PIP centres.
(3) Temporal trajectory plots: frame-by-frame main viewport centre position 
    for Mask R-CNN + KBRS vs.\ Director-CenterNet, illustrating reduced VD.}

%%====================================================================
\section{Discussion}
\label{sec:discussion}
%%====================================================================

\subsection{Why Multi-Region Finding Matters for Esports Broadcasting}

The fundamental tension in Esports broadcasting is that a single camera 
cannot cover the full extent of a large, concurrent game state. Real-time 
strategy games such as StarCraft exhibit up to 40 times the unit count of 
MOBA games \cite{joo2023learning}, making multi-region coverage not merely 
a convenience but a broadcast necessity. Director-CenterNet's dynamic PIP 
allocation based on KBRS residual salience addresses this need directly, 
without requiring additional human annotations beyond the single-viewport 
training data already used by prior work.

\subsection{CenterNet vs.\ Mask R-CNN for MRVP}

CenterNet offers three structural advantages over Mask R-CNN in the MRVP 
setting. First, its dense heatmap is directly commensurable with the KBRS 
score field, making coverage-guided PIP placement natural. Second, local-
maximum extraction is more stable and efficient than repeated non-maximum 
suppression over proposal boxes. Third, the Gaussian ground-truth rendering 
provides soft supervision that accommodates the inherent uncertainty in 
human viewport consensus, whereas Mask R-CNN's hard instance masks require 
a discrete assignment step.

\subsection{Toward Domain-Agnostic Multi-Region Observers}

The MRVP formulation and Director-CenterNet architecture are 
designed to be domain-agnostic. The structured game-state tensor can be 
replaced by any multi-channel spatial representation (e.g., MOBA minimap 
channels, drone swarm positional heatmaps, multi-CCTV occupancy grids). 
The KBRS scoring kernels can be adapted to domain-specific saliency notions 
(e.g., vehicle density plus pedestrian-vehicle mixture for traffic 
surveillance). The four-loss framework remains applicable as long as (i) 
human attention data is available for main-viewport supervision and (ii) a 
domain-relevant saliency score field can be defined for PIP placement.

\subsection{Limitations}

\paragraph{Track~B (zero human-GT setting).} Director-CenterNet currently 
operates in Track~A (human observation data available for the main viewport). 
Extending to Track~B — where no human viewport data exists and the system 
must rely entirely on unsupervised saliency signals — is a critical direction 
for cross-game generalisation and is left as future work.

\paragraph{Temporal modelling.} $\mathcal{L}_{\text{smooth}}$ acts on 
pairs of consecutive frames. Incorporating longer-range temporal context via 
recurrent modules or temporal self-attention may further reduce thrashing 
for events that evolve over several seconds.

\paragraph{PIP count.} We fix $K=3$ in all experiments. Adapting $K$ 
dynamically beyond the confidence-threshold mechanism could improve coverage 
when more than two concurrent events occur simultaneously.

\paragraph{Live integration.} All experiments are conducted in offline 
replay mode. Real-time integration requires addressing the pipeline latency 
between game-state extraction, inference, and camera actuation.

%%====================================================================
\section{Conclusion}
\label{sec:conclusion}
%%====================================================================

We presented Director-CenterNet, a multi-region viewport prediction 
framework for automated Esports broadcasting. By reformulating automatic 
observing as Multi-Region Viewport Prediction and introducing four 
complementary training losses — Human Consensus Match, KBRS Coverage, 
Spatial Repulsion, and Trajectory Smoothness — Director-CenterNet jointly 
predicts a main broadcast viewport and multiple PIP viewports from 
multi-channel game-state tensors in a single forward pass. The KBRS 
Coverage loss enables principled spatial resource allocation to salient 
uncovered regions without additional human annotations, while the Spatial 
Repulsion and Trajectory Smoothness losses eliminate redundant overlap and 
camera thrashing respectively. Experiments on StarCraft: Brood War benchmarks 
demonstrate improvements in viewport overlap, multi-viewport coverage, and 
temporal stability compared to all single-viewport baselines.

Future work will explore Track~B (zero human-GT) operation via unsupervised 
saliency estimation, longer-range temporal context modelling, and extension 
to MOBA games and non-game surveillance domains.

%%====================================================================
%% Declaration
%%====================================================================

\section*{CRediT authorship contribution statement}
\todo{Fill in author contributions following CRediT taxonomy.}

\section*{Declaration of competing interest}
The authors declare that they have no known competing financial interests 
or personal relationships that could have appeared to influence the work 
reported in this paper.

\section*{Acknowledgements}
\todo{Add funding acknowledgements.}

%%====================================================================
\appendix

\section{KBRS Kernel Definitions}
\label{app:kbrs_kernels}

For completeness we define the three KBRS kernels formally. Let 
$\mathbf{F}_t \in \mathbb{R}^{C \times H_s \times W_s}$ be the chosen FPN 
feature map. For a channel group $\mathcal{G} \subseteq \{1,\ldots,C\}$, 
define the grouped activation map
\begin{equation}
  G^{\mathcal{G}}_t(x,y) = \sum_{c \in \mathcal{G}} \mathbf{F}_t[c,x,y].
\end{equation}

\paragraph{Density (Eq.~\ref{eq:density}).}
$\mathbf{K}_1 \in \mathbb{R}^{k_h \times k_w}$ is an all-ones rectangular 
kernel. Convolution accumulates entity presence within a viewport-sized 
window.

\paragraph{Mixture (Eq.~\ref{eq:mixture}).}
Windowed sums $u_t = \mathbf{K}_1 * G^{\mathcal{G}_1}_t$ and 
$v_t = \mathbf{K}_1 * G^{\mathcal{G}_2}_t$ measure player-1 and player-2 
presence. The mixture score $S^{\text{mix}}_t = (4\rho_t(1-\rho_t))^\eta$ 
peaks when both players are equally present in a window (balanced 
co-presence), identifying likely combat or confrontation zones.

\paragraph{Centeredness (Eq.~\ref{eq:centeredness}).}
$\mathbf{K}_{\text{ctr}}$ is a 2D Gaussian kernel that assigns higher 
weight to channels near the kernel centre, encouraging viewports whose 
content is well-centred rather than clipped at the boundary.

\section{Hyperparameter Summary}
\label{app:hyperparams}

\begin{table}[h]
  \centering
  \caption{Hyperparameter settings used in all experiments unless 
  otherwise stated.}
  \label{tab:hyperparams}
  \begin{tabular}{lll}
    \toprule
    Hyperparameter & Value & Description \\
    \midrule
    Backbone         & ResNet-50  & Pretrained on ImageNet \\
    FPN stride $s$   & 32         & Coarsest FPN level \\
    Gaussian $\sigma$  & 2 tiles  & GT heatmap rendering \\
    Stacking $\Delta$  & 8 frames & Temporal stride \\
    Stacking $W$       & 3        & Number of past frames \\
    KBRS $\alpha$      & 0.3      & Density weight \\
    KBRS $\beta$       & 3.0      & Mixture weight \\
    KBRS $\gamma$      & 0.3      & Centeredness weight \\
    KBRS window        & $12{\times}20$ tiles & Viewport-matched \\
    $\lambda_{\text{off}}$   & 1.0 & Offset loss weight \\
    $\lambda_{\text{sz}}$    & 0.1 & Size loss weight \\
    $\lambda_{\text{cov}}$   & 0.5 & Coverage loss weight \\
    $\lambda_{\text{rep}}$   & 0.3 & Repulsion loss weight \\
    $\lambda_{\text{sm}}$    & 0.2 & Smoothness loss weight \\
    $K$              & 3          & Main + 2 PIP viewports \\
    $\tau$           & 0.2        & PIP confidence threshold \\
    Optimiser        & SGD        & momentum 0.9 \\
    Learning rate    & $5{\times}10^{-3}$ & StepLR (step 3, $\gamma$ 0.1) \\
    Weight decay     & $5{\times}10^{-4}$ & \\
    Batch size       & 16         & \\
    Epochs           & 30         & \\
    \bottomrule
  \end{tabular}
\end{table}

%%====================================================================
%% Bibliography
%%====================================================================

\bibliographystyle{elsarticle-num}
\bibliography{references}

\end{document}
